\section{The cost of choosing all queries in advance}
A nonadaptive classical tester chooses its full list before seeing the answers. In the quantum version, it prepares any entangled state, applies the bit oracle to all its query registers in one batch, and performs a measurement. Let $R^{\rm na}_\xi,Q^{\rm na}_\xi$ denote their bounded-error query costs.
\begin{theorem}\label{th:nonadaptive}
For fixed $\eta,\epsilon>0$, let $(1+\eta)n\le s\le(1-\epsilon)N/n$. For every $1/N\le\xi\le\xi_0$, where $\xi_0>0$ is a sufficiently small fixed constant depending on the margins,
\[
              Q^{\rm na}_\xi=\Theta_{\eta,\epsilon}(\min\{N,T/\xi\})
                   =R^{\rm na}_\xi .
\]
Both bounds also hold when the tester must accept every table within $\theta\xi N$ of $A_s$, for any fixed $0\le\theta<1$. The upper-bound constant may depend on $\theta$.
\end{theorem}
\begin{proof}
For the upper bound, query all entries or draw $O_\theta(T/\xi)$ uniform random entries. Accept if some member of $A_s$ has observed disagreement below $(1+\theta)\xi/2$. If a close member exists, its empirical disagreement obeys a multiplicative Chernoff upper-tail bound, using stochastic domination by sampling from a population with disagreement $\theta\xi$. On a far input every member has actual disagreement at least $\xi$. Its lower-tail failure has probability $\exp(-\Omega_\theta(\xi q))$ after $q$ samples, and (1) allows a union bound. Here observed disagreement denotes the sample fraction. Coincidences among sampled entries can be queried only once.

The lower bound uses a randomly chosen block of addresses. First suppose a power of two $M\in[4T,8T)$ fits inside $N$. Even if $M$ lies slightly above $N/2$, this argument applies. Use a uniform word in the binary linear space of Section~3 with length $M$ and small constant rate parameter $\rho$. For some $d=\lfloor\delta M\rfloor$ it is $d$-wise uniform. Every such word, interpreted as a Boolean table on $\log M$ address bits, is computable with at most $(9\rho+o(1))M/\log M$ gates. Take $\rho$ sufficiently small in terms of $\eta$, leaving at least $n+O(1)$ gates. In the hard experiment take $M$ independent fair bits. A fixed small $\alpha>0$ with $H_2(\alpha)<1/4$ ensures this word is at distance at least $\alpha M$ from every $s$-gate table on $\log M$ bits with probability $1-o(1)$. Indeed the circuit census on these bits has logarithm $(1+o(1))T\le(1+o(1))M/4$.

Let $p$ be the smallest reciprocal power of two at least $\max(M/N,\xi/\alpha)$ and choose $\xi_0\le\alpha/2$. Partition the table by a prefix of $\log(1/p)$ address bits. Choose one prefix uniformly. On its block write the chosen word on the last $\log M$ bits, ignoring any intervening bits. Set all other addresses to zero. The selector of a prefix of length $b$ uses at most $b+O(1)$ binary AND, OR and NOT gates: collect all required zeros with one OR followed by NOT, all required ones with one AND, then combine. The final gate combines this selector with the word, so the easy experiment stays within budget. In the hard experiment fix the chosen prefix and each assignment of the intervening bits. Every size-$s$ approximant on the full domain restricts to a size-$s$ approximant on this slice. Summation of its discrepancies gives distance at least $p\alpha N\ge\xi N$. The choice of $p$ gives $M/p=\Theta(\min(N,T/\xi))$.

Consider any one-batch $q$-query quantum tester, purifying its state $|\psi\rangle$ if necessary. For a fixed chosen block, let $\Pi$ keep the computational-basis query tuples with at most $\lfloor d/2\rfloor$ registers inside the block. Average over the uniform block. The expected number of active registers is $qp$, so $\mathbb E\|(I-\Pi)\psi\|^2\le2qp/d$. Replacing $|\psi\rangle$ by the unnormalized $\Pi|\psi\rangle$ changes the acceptance probability by at most $2\|(I-\Pi)\psi\|$, even after the oracle and the final measurement. Once projected, the average output density matrix is identical for the two word laws at this fixed block. Each matrix entry involves one ket query tuple and one bra query tuple, which together consult at most $d$ word coordinates. The two laws have the same joint marginal there. Hence the unprojected acceptance gap is at most $4\sqrt{2qp/d}$. A bounded-error tester on our promised inputs has a constant gap, implying $q=\Omega(d/p)$.

If no such $M$ fits, then $N<8T$. Use the code on all $N$ entries with $\rho$ small enough that every word costs at most $s$. By (1) and the fixed Shannon margin a uniform word is $\xi_0 N$-far with probability $1-o(1)$ when $H_2(\xi_0)<\epsilon/2$. Its marginals and the codeword's agree on every $\delta N$ coordinates, giving $\Omega(N)$ even adaptively by the polynomial method. The constructions use exact members on the easy side, and so the lower bound applies to every stated tolerance.
\end{proof}
In the exact total case a general theorem of Montanaro already forces linear nonadaptive quantum query complexity for every Boolean decision function depending on all $N$ entries \cite{mont}. The hidden block deals with a Hamming-distance promise, where neighboring opposite inputs need not exist. In the range shared with Theorem~\ref{th:testing}, the two results quantify the advantage of adaptive queries at shrinking distance.
